\documentclass[11pt]{article}
\usepackage[margin=1in]{geometry}
\usepackage[T1]{fontenc}
\usepackage[utf8]{inputenc}
\usepackage{enumitem}
\usepackage{xcolor}
\usepackage{longtable}
\usepackage{array}
\usepackage{hyperref}
\usepackage{listings}

\lstdefinestyle{code}{
  basicstyle=\ttfamily\small,
  columns=fullflexible,
  breaklines=true,
  frame=single,
  keepspaces=true
}

\newcommand{\file}[1]{\texttt{#1}}
\newcommand{\lineitem}[2]{\textbf{Line #1:} #2\\}

\title{UARTSim Code Explanation}
\author{Generated Documentation}
\date{\today}

\begin{document}
\maketitle

\section*{Overall Context}
UARTSim is a small console project that simulates a very simple UART transfer between two UART objects. The program is intentionally kept small:
\begin{itemize}[leftmargin=*]
  \item \file{src/main.cpp} drives the whole demo and user interaction.
  \item \file{include/Uart.hpp} defines the UART class, its registers, and its behavior.
  \item \file{include/Bus.hpp} defines a tiny bus that calls \texttt{tick()} on each UART.
  \item \file{include/CircularBuffer.hpp} defines a small reusable circular buffer template used by UART for TX/RX storage.
\end{itemize}

The data flow is:
\begin{enumerate}[leftmargin=*]
  \item The user enters UART options in the console menu.
  \item \texttt{main.cpp} creates two UART objects and a bus.
  \item One UART queues the message into its TX buffer.
  \item The bus repeatedly calls \texttt{tick()} on each UART.
  \item Each tick moves one byte from TX to the peer's RX buffer, or to itself in loopback mode.
  \item The program prints the received message and a register report.
\end{enumerate}

\section*{File: src/main.cpp}
\begin{lstlisting}[style=code]
#include "../include/Bus.hpp"
#include "../include/Uart.hpp"

#include <cstdint>
#include <iostream>
#include <sstream>
#include <string>

using namespace std;

namespace {
string readLineOrDefault(const string& prompt, const string& fallback) {
    cout << prompt << " [" << fallback << "]: ";
    string text;
    if (!getline(cin, text) || text.empty()) {
        return fallback;
    }
    return text;
}

int readIntOrDefault(const string& prompt, int fallback, int minValue, int maxValue) {
    while (true) {
        cout << prompt << " [" << fallback << "]: ";
        string text;
        if (!getline(cin, text) || text.empty()) {
            return fallback;
        }

        stringstream stream(text);
        int value = 0;
        char extra = '\0';
        if (stream >> value && !(stream >> extra) && value >= minValue && value <= maxValue) {
            return value;
        }

        cout << "Please enter a number from " << minValue << " to " << maxValue << ".\n";
    }
}

bool readYesNoOrDefault(const string& prompt, bool fallback) {
    const string fallbackText = fallback ? "y" : "n";
    while (true) {
        cout << prompt << " [" << fallbackText << "]: ";
        string text;
        if (!getline(cin, text) || text.empty()) {
            return fallback;
        }

        if (text == "y" || text == "Y" || text == "yes" || text == "YES") {
            return true;
        }
        if (text == "n" || text == "N" || text == "no" || text == "NO") {
            return false;
        }

        cout << "Please enter y or n.\n";
    }
}

Uart::ParityMode parityFromChoice(int choice) {
    if (choice == 1) {
        return Uart::ParityMode::Even;
    }
    if (choice == 2) {
        return Uart::ParityMode::Odd;
    }
    return Uart::ParityMode::None;
}

string parityName(Uart::ParityMode mode) {
    if (mode == Uart::ParityMode::Even) {
        return "even";
    }
    if (mode == Uart::ParityMode::Odd) {
        return "odd";
    }
    return "none";
}

string formatStatus(Uart::Registers registers) {
    stringstream out;
    out << "0x" << hex << registers.status << dec;
    out << " [TX_BUSY=" << ((registers.status & Uart::StatusTxBusy) != 0U ? "1" : "0")
        << ", RX_READY=" << ((registers.status & Uart::StatusRxReady) != 0U ? "1" : "0")
        << ", OVERFLOW=" << ((registers.status & Uart::StatusBufferOverflow) != 0U ? "1" : "0") << "]";
    return out.str();
}

void printStatusPhase(const string& label, const Uart& uart) {
    cout << label << " STATUS = " << formatStatus(uart.registers()) << '\n';
}

void printBuffer(Uart& uart) {
    uint8_t byte = 0U;
    while (uart.readReceivedByte(byte)) {
        cout << static_cast<char>(byte);
    }
}

void printRegisterReport(const Uart& uart, const string& receiverName) {
    const Uart::Registers registers = uart.registers();
    const Uart::Config config = uart.config();

    cout << "\nRegister report for " << uart.name() << "\n";
    cout << "CONTROL = 0x" << hex << registers.control << dec << '\n';
    cout << "  parity enable: " << (config.parityEnabled ? "set" : "clear")
         << " because parity was " << parityName(config.parity) << '\n';
    cout << "  parity odd: " << ((registers.control & Uart::ControlParityOdd) != 0U ? "set" : "clear")
         << " because " << parityName(config.parity) << " parity uses the odd bit only when needed" << '\n';
    cout << "  stop bits: " << static_cast<int>(config.stopBits)
         << " because the menu selected that frame length" << '\n';
    cout << "  data bits: " << static_cast<int>(config.dataBits)
         << " because the menu selected that word size" << '\n';
    cout << "  loopback: " << (config.loopback ? "set" : "clear")
         << " because bytes go " << (config.loopback ? "back to the same UART" : receiverName) << '\n';

    cout << "BAUD_DIV = " << registers.baudDiv << '\n';
    cout << "  why: derived from 16000000 / " << config.baudRate << " using the selected baud rate" << '\n';

    cout << "STATUS = 0x" << hex << registers.status << dec << '\n';
    cout << "  tx busy: " << ((registers.status & Uart::StatusTxBusy) != 0U ? "set" : "clear")
         << " because transmission is " << (uart.hasPendingTx() ? "still running" : "finished") << '\n';
    cout << "  rx ready: " << ((registers.status & Uart::StatusRxReady) != 0U ? "set" : "clear")
         << " because received bytes are " << (uart.hasPendingRx() ? "waiting in the buffer" : "already read") << '\n';
}
} // namespace

int main(int argc, char* argv[]) {
    string message = argc > 1 ? argv[1] : "Hello UART";

    cout << "UARTSim setup\n";
    cout << "Press Enter to accept each default.\n\n";

    message = readLineOrDefault("Message to send", message);
    const int parityChoice = readIntOrDefault("Parity 0=none 1=even 2=odd", 0, 0, 2);
    const int stopBitsChoice = readIntOrDefault("Stop bits 1 or 2", 1, 1, 2);
    const int dataBitsChoice = readIntOrDefault("Data bits 7 or 8", 8, 7, 8);
    const int baudRateChoice = readIntOrDefault("Baud rate", 9600, 300, 1000000);
    const bool loopbackChoice = readYesNoOrDefault("Loopback mode y/n", false);

    Uart uart0{"UART0"};
    Uart uart1{"UART1"};
    Bus bus;

    Uart::Config config;
    config.parity = parityFromChoice(parityChoice);
    config.parityEnabled = config.parity != Uart::ParityMode::None;
    config.stopBits = static_cast<std::uint8_t>(stopBitsChoice);
    config.dataBits = static_cast<std::uint8_t>(dataBitsChoice);
    config.baudRate = static_cast<std::uint32_t>(baudRateChoice);
    config.loopback = loopbackChoice;

    uart0.configure(config);
    uart1.configure(config);

    uart0.connectPeer(&uart1);
    uart1.connectPeer(&uart0);
    bus.attach(uart0);
    bus.attach(uart1);

    uart0.send(message);

    cout << "UART demo\n";
    cout << "Sent: " << message << '\n';
    printStatusPhase("At transmission time:", uart0);

    int cycle = 0;
    while (uart0.hasPendingTx()) {
        ++cycle;
        bus.tick();
        printStatusPhase("During transmission cycle " + to_string(cycle) + ":", uart0);
    }

    Uart& receiver = config.loopback ? uart0 : uart1;
    const string receiverName = receiver.name();

    printStatusPhase("After transmission:", uart0);

    cout << "\nReceived on " << receiverName << ": ";
    printBuffer(receiver);
    cout << '\n';

    printRegisterReport(uart0, receiverName);
    return 0;
}
\end{lstlisting}

\subsection*{Line-by-line explanation}
\begin{longtable}{>{\raggedright\arraybackslash}p{0.12\textwidth} p{0.82\textwidth}}
\textbf{Line} & \textbf{Explanation} \\
\hline
\endhead
1 & Includes \file{Bus.hpp} using a relative path so the compiler can find it even if the build is started from \file{src/}. \\
2 & Includes \file{Uart.hpp} using a relative path for the same reason. \\
3--6 & Standard library headers for fixed-width integers, console I/O, string parsing, and text handling. \\
8 & Brings the \texttt{std} names into the global scope to keep the code shorter. \\
10 & Starts an anonymous namespace so helper functions stay local to this file. \\
11--17 & \texttt{readLineOrDefault} asks for a text value and returns the fallback if the user presses Enter. \\
12 & Prints the prompt and the default value. \\
13--16 & Reads a full line from the console and returns the fallback if the line is empty. \\
19--33 & \texttt{readIntOrDefault} asks for an integer inside a valid range. \\
20--21 & Keeps asking until a valid number is entered. \\
22--25 & Reads the text and uses a \texttt{stringstream} to parse it as an integer. \\
27--29 & Verifies that the parse succeeded, no extra characters remain, and the value is inside the allowed range. \\
31 & Prints an error message if the input is invalid. \\
36--52 & \texttt{readYesNoOrDefault} reads a yes/no choice and accepts \texttt{y}, \texttt{n}, \texttt{yes}, or \texttt{no}. \\
37 & Converts the fallback boolean into a display string. \\
38--50 & Repeats until the user enters a valid yes/no value or presses Enter for default. \\
54--61 & Converts a menu choice into a UART parity mode. \\
55--57 & Choice 1 means even parity. \\
58--60 & Choice 2 means odd parity. \\
61 & Anything else means parity is disabled. \\
64--72 & Converts a parity mode back into a printable name. \\
74--80 & Formats the status register as hex plus named bits for easier reading. \\
75--79 & Builds a string showing \texttt{TX\_BUSY}, \texttt{RX\_READY}, and \texttt{OVERFLOW}. \\
83--85 & Prints a single status checkpoint with a label. \\
87--92 & Reads bytes from the receive buffer and prints them as characters. \\
94--117 & Prints a human-readable report of the UART registers and explains why each field has the value it has. \\
95--96 & Copies the current register and config state so the report can inspect them. \\
98 & Starts the per-register explanation block for \texttt{CONTROL}. \\
99--109 & Explains parity, stop bits, data bits, loopback, and baud divisor. \\
111--116 & Explains the status register and why \texttt{TX\_BUSY} and \texttt{RX\_READY} are set or clear. \\
118 & Ends the anonymous namespace. \\
120--155 & \texttt{main} performs the actual demo. \\
121 & Uses the first command-line argument as the message, or defaults to \texttt{Hello UART}. \\
123--124 & Prints the setup banner. \\
126--130 & Collects the message, parity, stop bits, data bits, baud rate, and loopback option from the user. \\
132--134 & Creates two UART objects and one bus object. \\
136--142 & Fills the UART configuration struct from the menu choices. \\
144--145 & Applies the configuration to both UARTs. \\
147--149 & Connects the UARTs as peers and attaches them to the bus. \\
151 & Queues the message for transmission. \\
153--154 & Prints the initial transmission status before the bus starts ticking. \\
156--160 & Repeatedly ticks the bus until the transmit buffer becomes empty, printing the status after each cycle. \\
162--163 & Chooses which UART should be treated as the receiver, depending on loopback mode. \\
165--166 & Prints the final status after transmission is complete. \\
168--170 & Prints the received text. \\
172--173 & Prints the detailed register report. \\
174 & Returns success. \\
\end{longtable}

\section*{File: include/Uart.hpp}
\begin{lstlisting}[style=code]
#pragma once

#include "CircularBuffer.hpp"

#include <cstdint>
#include <string>
#include <utility>
using namespace std;

class Uart {
public:
    enum class ParityMode {
        None,
        Even,
        Odd,
    };

    struct Config {
        ParityMode parity = ParityMode::None;
        bool parityEnabled = false;
        uint8_t stopBits = 1U;
        uint8_t dataBits = 8U;
        uint32_t baudRate = 9600U;
        bool loopback = false;
    };

    struct Registers {
        uint32_t control = 0U;
        uint32_t baudDiv = 1U;
        uint32_t status = 0U;
    };

    enum ControlBit : std::uint32_t {
        ControlParityEnable = 1U << 0,
        ControlParityOdd = 1U << 1,
        ControlTwoStopBits = 1U << 2,
        ControlSevenDataBits = 1U << 3,
        ControlLoopback = 1U << 4,
    };

    enum StatusBit : std::uint32_t {
        StatusTxBusy = 1U << 0,
        StatusRxReady = 1U << 1,
        StatusBufferOverflow = 1U << 2,
    };

    explicit Uart(std::string name)
        : name_(std::move(name)) {
    }

    void connectPeer(Uart* peer) {
        peer_ = peer;
    }

    void configure(const Config& config) {
        config_ = config;
        if (!config_.parityEnabled) {
            config_.parity = ParityMode::None;
        }

        registers_.control = 0U;
        if (config_.parityEnabled) {
            registers_.control |= ControlParityEnable;
            if (config_.parity == ParityMode::Odd) {
                registers_.control |= ControlParityOdd;
            }
        }
        if (config_.stopBits == 2U) {
            registers_.control |= ControlTwoStopBits;
        }
        if (config_.dataBits == 7U) {
            registers_.control |= ControlSevenDataBits;
        }
        if (config_.loopback) {
            registers_.control |= ControlLoopback;
        }

        registers_.baudDiv = config_.baudRate == 0U ? 1U : (BaseClockHz / config_.baudRate);
        if (registers_.baudDiv == 0U) {
            registers_.baudDiv = 1U;
        }

        registers_.status = 0U;
    }

    void send(const std::string& text) {
        for (char ch : text) {
            if (!txBuffer_.push(static_cast<std::uint8_t>(ch))) {
                registers_.status |= StatusBufferOverflow;
                break;
            }
        }

        if (!txBuffer_.empty()) {
            registers_.status |= StatusTxBusy;
        }
    }

    void tick() {
        uint8_t byte = 0U;
        if (!txBuffer_.pop(byte)) {
            if (txBuffer_.empty()) {
                registers_.status &= static_cast<std::uint32_t>(~StatusTxBusy);
            }
            return;
        }

        if (config_.loopback) {
            if (rxBuffer_.push(byte)) {
                registers_.status |= StatusRxReady;
            }
        } else if (peer_ != nullptr) {
            peer_->rxBuffer_.push(byte);
            peer_->registers_.status |= StatusRxReady;
        }

        if (txBuffer_.empty()) {
            registers_.status &= static_cast<std::uint32_t>(~StatusTxBusy);
        }
    }

    bool readReceivedByte(uint8_t& byte) {
        const bool popped = rxBuffer_.pop(byte);
        if (rxBuffer_.empty()) {
            registers_.status &= static_cast<uint32_t>(~StatusRxReady);
        }
        return popped;
    }

    bool hasPendingTx() const {
        return !txBuffer_.empty();
    }

    bool hasPendingRx() const {
        return !rxBuffer_.empty();
    }

    const string& name() const {
        return name_;
    }

    const Config& config() const {
        return config_;
    }

    const Registers& registers() const {
        return registers_;
    }

private:
    static constexpr std::uint32_t BaseClockHz = 16000000U;

    string name_;
    Uart* peer_ = nullptr;
    Config config_{};
    Registers registers_{};
    CircularBuffer<std::uint8_t, 16U> txBuffer_;
    CircularBuffer<std::uint8_t, 16U> rxBuffer_;
};
\end{lstlisting}

\subsection*{Line-by-line explanation}
\begin{longtable}{>{\raggedright\arraybackslash}p{0.12\textwidth} p{0.82\textwidth}}
\textbf{Line} & \textbf{Explanation} \\
\hline
\endhead
1 & Prevents multiple inclusion of this header. \\
3 & Includes the circular buffer template used for transmit and receive storage. \\
5--6 & Standard headers for fixed-width integers and strings. \\
7 & Imports \texttt{std} symbols into the header. This works, but in general it is better not to do this in headers. \\
9 & Starts the \texttt{Uart} class. \\
10--20 & Defines the parity mode enum. \\
22--28 & Defines a configuration object that stores user-selected UART options. \\
30--34 & Defines a register bundle representing control, baud divisor, and status. \\
36--43 & Defines bit flags for the control register. \\
45--50 & Defines bit flags for the status register. \\
52--54 & Constructor stores the UART name. \\
56--58 & Connects this UART to its peer. \\
60--81 & Applies the selected configuration and sets the virtual registers accordingly. \\
61 & Copies the config into the object. \\
62--64 & Forces parity off if the user disabled it. \\
66--77 & Builds the control register by setting bits for parity, stop bits, data bits, and loopback. \\
79--82 & Computes the baud divisor from the fixed base clock. \\
84 & Clears status bits when configuration is applied. \\
86--95 & Queues a string into the TX buffer one byte at a time. \\
88--90 & If the TX buffer is full, marks overflow and stops queuing more bytes. \\
92--94 & Marks the UART busy when there is pending TX data. \\
96--114 & One tick transfers one byte. \\
97--100 & Pops one byte from TX; if nothing is pending, clears busy and returns. \\
102--106 & If loopback is enabled, the byte goes back into the same UART's RX buffer. \\
107--110 & Otherwise the byte is pushed into the peer UART's RX buffer and the peer RX-ready bit is set. \\
112--114 & Clears TX busy when the queue becomes empty. \\
116--120 & Reads one received byte and clears RX-ready when the RX buffer becomes empty. \\
122--130 & Status queries and name accessor functions. \\
132--137 & Accessors for config and registers. \\
139--147 & Private data members holding the name, peer pointer, config, registers, and buffers. \\
140 & Base clock used to derive baud divisor. \\
\end{longtable}

\section*{File: include/Bus.hpp}
\begin{lstlisting}[style=code]
#pragma once

#include "Uart.hpp"

#include <vector>

class Bus {
public:
    void attach(Uart& uart) {
        uarts_.push_back(&uart);
    }

    void tick() {
        for (Uart* uart : uarts_) {
            uart->tick();
        }
    }

private:
    std::vector<Uart*> uarts_;
};
\end{lstlisting}

\subsection*{Line-by-line explanation}
\begin{longtable}{>{\raggedright\arraybackslash}p{0.12\textwidth} p{0.82\textwidth}}
\textbf{Line} & \textbf{Explanation} \\
\hline
\endhead
1 & Prevents multiple inclusion. \\
3 & Includes the UART class definition because the bus stores UART pointers. \\
5 & Includes \texttt{vector} so the bus can hold a dynamic list of UARTs. \\
7 & Declares the bus class. \\
8--10 & \texttt{attach} stores a UART pointer in the internal list. \\
12--16 & \texttt{tick} loops over all attached UARTs and calls their \texttt{tick()} method. \\
18--19 & The bus stores pointers to UART objects. \\
\end{longtable}

\section*{File: include/CircularBuffer.hpp}
\begin{lstlisting}[style=code]
#pragma once

#include <array>
#include <cstddef>
using namespace std;
template <typename T, std::size_t Capacity>
class CircularBuffer {
public:
    bool push(const T& value) {
        if (count_ == Capacity) {
            return false;
        }

        values_[tail_] = value;
        tail_ = (tail_ + 1U) % Capacity;
        ++count_;
        return true;
    }

    bool pop(T& value) {
        if (count_ == 0U) {
            return false;
        }

        value = values_[head_];
        head_ = (head_ + 1U) % Capacity;
        --count_;
        return true;
    }

    bool empty() const {
        return count_ == 0U;
    }

    bool full() const {
        return count_ == Capacity;
    }

private:
    array<T, Capacity> values_{};
    size_t head_ = 0U;
    size_t tail_ = 0U;
    size_t count_ = 0U;
};
\end{lstlisting}

\subsection*{Line-by-line explanation}
\begin{longtable}{>{\raggedright\arraybackslash}p{0.12\textwidth} p{0.82\textwidth}}
\textbf{Line} & \textbf{Explanation} \\
\hline
\endhead
1 & Prevents multiple inclusion. \\
3--4 & Includes \texttt{array} for fixed-size storage and \texttt{cstddef} for \texttt{size\_t}. \\
5 & Imports \texttt{std} symbols into the header. This is convenient but not ideal in headers. \\
6 & Begins the template declaration for a circular buffer storing type \texttt{T} with capacity \texttt{Capacity}. \\
7 & Starts the class definition. \\
8--17 & \texttt{push} inserts a value if the buffer is not full. \\
9--11 & Returns false when the buffer is already full. \\
13--15 & Stores the new value, advances the tail index, and increments the element count. \\
18--27 & \texttt{pop} removes a value if the buffer is not empty. \\
19--21 & Returns false when the buffer is empty. \\
23--25 & Copies out the front value, advances the head index, and decrements the count. \\
28--30 & Returns whether the buffer currently holds zero elements. \\
32--34 & Returns whether the buffer is completely full. \\
36--40 & Internal storage and state variables for the ring buffer. \\
37 & Fixed-size array holding the actual items. \\
38--40 & Indices and count used to manage wraparound behavior. \\
\end{longtable}

\section*{File: README.md}
The README explains the project at a high level. It says the program is a small C++17 console UART simulator with a text-based setup menu, shows how to build it with \file{build.bat}, and says to run \file{run.bat}. It also tells the user to enter message text, parity, stop bits, data bits, baud rate, and loopback in the menu.

\section*{File: CMakeLists.txt}
The CMake file defines a C++17 project named \texttt{UARTSim}, builds one executable from \file{src/main.cpp}, adds \file{include/} to the include path, and enables warnings. That means the project is intentionally simple: the logic lives in headers, and the program entry point is in one source file.

\section*{How the files work together}
\begin{itemize}[leftmargin=*]
  \item \file{main.cpp} is the driver. It reads user input, configures UART objects, starts transmission, and prints output.
  \item \file{Uart.hpp} is the core model. It stores configuration, register values, and the TX/RX buffers.
  \item \file{Bus.hpp} is the scheduler. It calls \texttt{tick()} repeatedly to simulate time passing.
  \item \file{CircularBuffer.hpp} is the utility container that makes the TX and RX queues work.
\end{itemize}

\section*{Important note}
The current headers contain \texttt{using namespace std;} inside the header files. The code still compiles, but in real projects it is better to avoid that in headers because it can leak names into other files.

\end{document}
